% TODO checklist:
% - [x] README.md: Feature boundaries, technology choices
% - [x] Q&A.md: Q46 (Limitations), Q50 (Comparison showing gaps)
% - [x] sapthesis-doc.tex: Non-objectives and scope notes in requirements chapter

\section{Scope and Non-Goals}
\label{sec:scope-nongoals}

Clearly delineating what falls within and outside the scope of this thesis is essential for understanding its contributions and limitations. This section explicitly defines the boundaries of the work.

\subsection{In-Scope Elements}

The following elements are within the scope of this thesis and are addressed through design, implementation, and evaluation:

\subsubsection{Agent Orchestration Engine}

\begin{itemize}
    \item Multi-step agent run execution with planning and tool invocation
    \item Intent classification pipeline (pattern-based and LLM fallback)
    \item Session management for conversational context
    \item Agent run lifecycle management (pending → running → completed/failed/cancelled)
    \item Step-by-step execution recording for auditability
    \item TODO list planning for complex task decomposition
\end{itemize}

\subsubsection{MCP Tool Framework}

\begin{itemize}
    \item Tool registry with 34 tools across 17 categories
    \item Tool discovery API with schema retrieval
    \item Tool invocation with input validation and output handling
    \item Per-tool RBAC and audit logging
    \item Patterns for tool extension and customization
\end{itemize}

\subsubsection{Graph Database Integration}

\begin{itemize}
    \item Memgraph graph database integration
    \item NL-to-Cypher pipeline for natural language querying
    \item Multi-layer safety validation for generated queries
    \item Graph schema management and introspection
    \item Result summarization for user-friendly responses
\end{itemize}

\subsubsection{Security Framework}

\begin{itemize}
    \item OIDC/JWT authentication with JWKS key management
    \item Role-Based Access Control with scope-based permissions
    \item Multi-tenant data isolation at database level
    \item Rate limiting with sliding window algorithm
    \item PII detection and redaction
    \item Comprehensive audit logging
\end{itemize}

\subsubsection{Observability Stack}

\begin{itemize}
    \item Prometheus metrics for all major subsystems
    \item Structured logging with correlation IDs
    \item OpenTelemetry distributed tracing
    \item Kubernetes-compatible health probes
    \item Component health monitoring
\end{itemize}

\subsubsection{Background Job System}

\begin{itemize}
    \item Asynchronous job queue with Redis backend
    \item Worker processes with heartbeat monitoring
    \item SSE-based progress streaming
    \item Job cancellation and graceful shutdown
    \item Scheduled background tasks (APScheduler)
\end{itemize}

\subsubsection{User Interfaces}

\begin{itemize}
    \item Next.js Agent Chat UI for end-user interactions
    \item Streamlit Control Panel for administration
    \item API client implementations in TypeScript and Python
\end{itemize}

\subsubsection{Deployment Infrastructure}

\begin{itemize}
    \item Docker Compose configuration for development and production
    \item NGINX reverse proxy configuration
    \item Prometheus and Grafana monitoring setup
    \item Database initialization and migration scripts
\end{itemize}

\subsection{Out-of-Scope Elements (Non-Goals)}

The following elements are explicitly \textbf{not} within the scope of this thesis:

\subsubsection{Advanced RAG Pipeline}

While the platform supports graph-based knowledge querying through NL-to-Cypher, it does \textbf{not} implement:

\begin{itemize}
    \item Document ingestion and chunking pipelines
    \item Vector embedding generation and storage
    \item Semantic similarity search over unstructured documents
    \item Hybrid retrieval combining vector and keyword search
\end{itemize}

\paragraph{Rationale:} The focus is on structured knowledge in graph form, which is the primary data representation at CINECA. Document RAG is a complementary capability that could be added in future work but is not essential for the initial deployment context.

\subsubsection{Multi-Agent Collaboration}

The platform implements single-agent orchestration but does \textbf{not} include:

\begin{itemize}
    \item Agent-to-agent communication protocols
    \item Multi-agent task delegation and coordination
    \item Hierarchical agent architectures (manager/worker patterns)
    \item Agent negotiation and conflict resolution
\end{itemize}

\paragraph{Rationale:} Multi-agent systems introduce significant complexity in coordination, debugging, and security. The single-agent architecture sufficiently addresses current requirements while providing a foundation for future multi-agent extensions.

\subsubsection{Advanced Agent Architectures}

The platform implements the ReAct (Reasoning + Acting) pattern but does \textbf{not} include:

\begin{itemize}
    \item Plan-and-execute architectures with separate planning and execution phases
    \item Tree-of-thought reasoning with branching exploration
    \item Self-reflection and critique loops
    \item Constitutional AI safety patterns
\end{itemize}

\paragraph{Rationale:} The ReAct pattern provides a solid foundation for tool-using agents. More sophisticated architectures could be implemented within the existing framework but are not required for the current use cases.

\subsubsection{Real-Time Streaming Responses}

The platform supports Server-Sent Events (SSE) for job progress but does \textbf{not} implement:

\begin{itemize}
    \item Token-by-token LLM response streaming to users
    \item WebSocket-based bidirectional communication
    \item Real-time collaborative editing features
\end{itemize}

\paragraph{Rationale:} The current polling-based approach provides adequate user experience for the target use cases. Streaming would improve perceived responsiveness but adds implementation complexity that was deprioritized.

\subsubsection{Code Execution Sandbox}

The platform does \textbf{not} include:

\begin{itemize}
    \item Sandboxed code execution environments for agent-generated code
    \item Container-based isolation for untrusted code
    \item Language-specific interpreters for code execution tools
\end{itemize}

\paragraph{Rationale:} Safe code execution requires sophisticated sandboxing that is outside the scope of this work. Agents can invoke pre-defined tools but cannot execute arbitrary code.

\subsubsection{Conversation Memory Systems}

Beyond session-scoped context, the platform does \textbf{not} implement:

\begin{itemize}
    \item Long-term episodic memory across sessions
    \item Semantic memory with embedding-based retrieval
    \item Summary memory for conversation compression
    \item Entity memory for tracking mentioned entities over time
\end{itemize}

\paragraph{Rationale:} Session-based context suffices for the current use cases. Advanced memory systems could enhance user experience but require additional infrastructure (vector stores) and careful privacy considerations.

\subsubsection{MLOps and Model Training}

The platform does \textbf{not} include:

\begin{itemize}
    \item Model fine-tuning or training pipelines
    \item Experiment tracking and model versioning
    \item A/B testing infrastructure for model comparison
    \item Model deployment automation
\end{itemize}

\paragraph{Rationale:} The platform consumes pre-trained models via APIs. Model training and MLOps are separate concerns addressed by dedicated platforms (MLflow, Kubeflow, etc.).

\subsection{Boundary Cases}

Some capabilities exist in limited form within the platform:

\begin{description}
    \item[Caching:] LLM response caching is supported but not extensively optimized. Semantic deduplication of similar prompts is not implemented.
    
    \item[Distributed Deployment:] While the architecture supports horizontal scaling, distributed locking for multi-worker job processing is not fully implemented.
    
    \item[Provider-Specific Features:] The platform uses a common provider interface, which means provider-specific features (e.g., OpenAI function calling syntax) are abstracted rather than directly exposed.
    
    \item[Graph Analytics:] Basic graph analytics (centrality, paths) are available through tools, but advanced algorithms (community detection, link prediction) are not implemented.
\end{description}

\subsection{Scope Summary}

Table~\ref{tab:scope-summary} provides a concise summary of scope decisions.

\begin{table}[ht]
\centering
\caption{Scope summary: included and excluded capabilities.}
\label{tab:scope-summary}
\small
\begin{tabular}{lcc}
\hline
\textbf{Capability} & \textbf{In Scope} & \textbf{Out of Scope} \\
\hline
Single-agent orchestration & \checkmark & \\
Multi-agent collaboration & & \checkmark \\
Graph NL-to-Cypher & \checkmark & \\
Document RAG & & \checkmark \\
OIDC/JWT authentication & \checkmark & \\
RBAC authorization & \checkmark & \\
Multi-tenancy & \checkmark & \\
Prometheus metrics & \checkmark & \\
OpenTelemetry tracing & \checkmark & \\
SSE job streaming & \checkmark & \\
Token streaming & & \checkmark \\
MCP tools (34) & \checkmark & \\
Code execution sandbox & & \checkmark \\
Session context & \checkmark & \\
Long-term memory & & \checkmark \\
Docker deployment & \checkmark & \\
Kubernetes operators & & \checkmark \\
\hline
\end{tabular}
\end{table}

\subsection{Extensibility for Future Work}

While certain capabilities are out of scope for this thesis, the architecture has been designed to accommodate future extensions:

\begin{itemize}
    \item The tool framework can incorporate document RAG tools
    \item The orchestrator can be extended with additional agent patterns
    \item The provider abstraction can support streaming providers
    \item The job system can be enhanced with distributed locking
    \item Memory systems can be added as additional tools or services
\end{itemize}

Chapter~\ref{ch:conclusions} discusses these extension opportunities in the context of future work.

